\textbf{Phần 1. Trắc nghiệm lựa chọn}
\setcounterpageref{ex}{0}

\begin{ex} %[5H9N51]
Công thức tính diện tích xung quanh của hình lập phương có cạnh $a$ là
\choice
{$S_{xq} = a \times a \times 6$}
{\True $S_{xq} = a \times a \times 4$}
{$S_{xq} = a \times 4$}
{$S_{xq} = a \times a \times a$}
\loigiai{
Diện tích xung quanh của hình lập phương bằng diện tích một mặt nhân với $4$: $S_{xq} = a \times a \times 4$.\\
\textbf{Đáp án của câu hỏi là B}
}
\end{ex}

\begin{ex} %[5H9N51]
Công thức tính diện tích toàn phần của hình lập phương có cạnh $a$ là
\choice
{\True $S_{tp} = a \times a \times 6$}
{$S_{tp} = a \times a \times 4$}
{$S_{tp} = a \times 6$}
{$S_{tp} = a \times a \times a$}
\loigiai{
Diện tích toàn phần của hình lập phương bằng diện tích một mặt nhân với $6$: $S_{tp} = a \times a \times 6$.\\
\textbf{Đáp án của câu hỏi là A}
}
\end{ex}

\begin{ex} %[5H9H51]
Một hình lập phương có cạnh $5\text{ cm}$. Diện tích xung quanh của hình lập phương đó là
\choice
{$25\text{ cm}^2$}
{$150\text{ cm}^2$}
{$125\text{ cm}^2$}
{\True $100\text{ cm}^2$}
\loigiai{
$S_{xq} = 5 \times 5 \times 4 = 100\text{ (cm}^2\text{)}$.\\
\textbf{Đáp án của câu hỏi là D}
}
\end{ex}

\begin{ex} %[5H9H51]
Một hình lập phương có cạnh $4\text{ dm}$. Diện tích toàn phần của hình lập phương đó là
\choice
{$64\text{ dm}^2$}
{$16\text{ dm}^2$}
{\True $96\text{ dm}^2$}
{$48\text{ dm}^2$}
\loigiai{
$S_{tp} = 4 \times 4 \times 6 = 96\text{ (dm}^2\text{)}$.\\
\textbf{Đáp án của câu hỏi là C}
}
\end{ex}

\begin{ex} %[5H9H51]
Một cái hộp hình lập phương không có nắp có cạnh $6\text{ cm}$. Diện tích bìa để làm chiếc hộp đó là
\choice
{$144\text{ cm}^2$}
{$216\text{ cm}^2$}
{$36\text{ cm}^2$}
{\True $180\text{ cm}^2$}
\loigiai{
Hộp không nắp có 5 mặt: $6 \times 6 \times 5 = 180\text{ (cm}^2\text{)}$.\\
\textbf{Đáp án của câu hỏi là D}
}
\end{ex}

\begin{ex} %[5H9H51]
Một hình lập phương có diện tích một mặt là $36\text{ cm}^2$. Diện tích toàn phần của hình lập phương đó là
\choice
{$144\text{ cm}^2$}
{$180\text{ cm}^2$}
{\True $216\text{ cm}^2$}
{$240\text{ cm}^2$}
\loigiai{
$S_{tp} = 36 \times 6 = 216\text{ (cm}^2\text{)}$.\\
\textbf{Đáp án của câu hỏi là C}
}
\end{ex}

\begin{ex} %[5H9V51]
Một hình lập phương có diện tích xung quanh là $144\text{ dm}^2$. Cạnh của hình lập phương đó là
\choice
{$4\text{ dm}$}
{$8\text{ dm}$}
{$12\text{ dm}$}
{\True $6\text{ dm}$}
\loigiai{
Diện tích một mặt: $144 : 4 = 36\text{ (dm}^2\text{)}$.\\
Vì $6 \times 6 = 36$ nên cạnh là $6\text{ dm}$.\\
\textbf{Đáp án của câu hỏi là D}
}
\end{ex}

\begin{ex} %[5H9V51]
Người ta sơn toàn bộ mặt ngoài của một khối gỗ hình lập phương có cạnh $0{,}8\text{ m}$. Diện tích cần sơn là
\choice
{$2{,}56\text{ m}^2$}
{$1{,}92\text{ m}^2$}
{\True $3{,}84\text{ m}^2$}
{$0{,}512\text{ m}^2$}
\loigiai{
$S_{tp} = 0{,}8 \times 0{,}8 \times 6 = 3{,}84\text{ (m}^2\text{)}$.\\
\textbf{Đáp án của câu hỏi là C}
}
\end{ex}

\begin{ex} %[5H9V51]
Một hình lập phương có diện tích toàn phần là $294\text{ cm}^2$. Diện tích xung quanh của hình lập phương đó là
\choice
{\True $196\text{ cm}^2$}
{$49\text{ cm}^2$}
{$98\text{ cm}^2$}
{$147\text{ cm}^2$}
\loigiai{
Diện tích một mặt: $294 : 6 = 49\text{ (cm}^2\text{)}$.\\
$S_{xq} = 49 \times 4 = 196\text{ (cm}^2\text{)}$.\\
\textbf{Đáp án của câu hỏi là A}
}
\end{ex}

\begin{ex} %[5H9V51]
Nếu gấp cạnh của một hình lập phương lên $3$ lần thì diện tích xung quanh của nó tăng lên gấp mấy lần?
\choice
{$3$ lần}
{\True $9$ lần}
{$6$ lần}
{$27$ lần}
\loigiai{
Diện tích một mặt tăng $3 \times 3 = 9$ lần nên diện tích xung quanh tăng lên 9 lần.\\
\textbf{Đáp án của câu hỏi là B}
}
\end{ex}

\begin{ex} %[5H9V51]
Người ta dán giấy màu lên tất cả các mặt của một chiếc hộp quà hình lập phương cạnh $12\text{ cm}$. Biết giá tiền giấy màu là $500$ đồng/$100\text{ cm}^2$. Tổng số tiền giấy màu để bọc hộp quà là
\choice
{\True $4\,320$ đồng}
{$3\,600$ đồng}
{$4\,800$ đồng}
{$5\,200$ đồng}
\loigiai{
$S_{tp} = 12 \times 12 \times 6 = 864\text{ (cm}^2\text{)}$.\\
Tiền giấy: $864 : 100 \times 500 = 4\,320$ (đồng).\\
\textbf{Đáp án của câu hỏi là A}
}
\end{ex}

\begin{ex} %[5H9C51]
Hai hình lập phương có tỉ số độ dài cạnh là $\dfrac{2}{3}$. Tỉ số diện tích toàn phần của hình thứ nhất so với hình thứ hai là
\choice
{$\dfrac{2}{3}$}
{\True $\dfrac{4}{9}$}
{$\dfrac{8}{27}$}
{$\dfrac{4}{6}$}
\loigiai{
Tỉ số diện tích toàn phần bằng bình phương tỉ số cạnh: $\left(\dfrac{2}{3}\right)^2 = \dfrac{4}{9}$.\\
\textbf{Đáp án của câu hỏi là B}
}
\end{ex}

\textbf{Phần 2. Trắc nghiệm đúng sai}
\setcounterpageref{ex}{0}

\begin{ex} %[5H9H51]
Cho một hình lập phương có cạnh bằng $8\text{ cm}$. Xét tính đúng sai:
\choiceTF
{\True Diện tích một mặt của hình lập phương là $64\text{ cm}^2$.}
{\True Diện tích xung quanh của hình lập phương là $256\text{ cm}^2$.}
{Diện tích toàn phần của hình lập phương là $320\text{ cm}^2$.}
{\True Diện tích toàn phần của hình lập phương là $384\text{ cm}^2$.}
\loigiai{
\begin{itemize}
    \item $S_{1mặt} = 8 \times 8 = 64\text{ cm}^2$. Đúng.
    \item $S_{xq} = 64 \times 4 = 256\text{ cm}^2$. Đúng.
    \item $S_{tp} = 64 \times 6 = 384\text{ cm}^2 \neq 320\text{ cm}^2$. Ý 3 Sai, Ý 4 Đúng.
\end{itemize}
\begin{itemize}
    \item a) \textbf{Đúng}.
    \item b) \textbf{Đúng}.
    \item c) \textbf{Sai}.
    \item d) \textbf{Đúng}.
\end{itemize}
}
\end{ex}

\begin{ex} %[5H9V51]
Một cái thùng hình lập phương không nắp có diện tích toàn bộ các mặt ngoài là $180\text{ dm}^2$. Xét tính đúng sai:
\choiceTF
{\True Chiếc thùng có $5$ mặt ngoài.}
{\True Diện tích một mặt của chiếc thùng là $36\text{ dm}^2$.}
{Cạnh của chiếc thùng là $8\text{ dm}$.}
{\True Diện tích xung quanh của chiếc thùng là $144\text{ dm}^2$.}
\loigiai{
\begin{itemize}
    \item Thùng không nắp có 5 mặt. Đúng.
    \item Diện tích 1 mặt: $180 : 5 = 36\text{ dm}^2$. Đúng.
    \item Cạnh thùng: $6\text{ dm}$ vì $6 \times 6 = 36$ (không phải 8 dm). Sai.
    \item $S_{xq} = 36 \times 4 = 144\text{ dm}^2$. Đúng.
\end{itemize}
\begin{itemize}
    \item a) \textbf{Đúng}.
    \item b) \textbf{Đúng}.
    \item c) \textbf{Sai}.
    \item d) \textbf{Đúng}.
\end{itemize}
}
\end{ex}

\textbf{Phần 3. Trả lời ngắn}
\setcounterpageref{ex}{0}

\begin{ex} %[5H9H51]
Một hình lập phương có cạnh $7\text{ cm}$. Tính diện tích xung quanh theo xăng-ti-mét vuông.
\shortans{196}
\loigiai{
\textbf{Đáp số:} $196$\\
$7 \times 7 \times 4 = 196\text{ cm}^2$.
}
\end{ex}

\begin{ex} %[5H9H51]
Tính diện tích toàn phần của hình lập phương có cạnh $9\text{ dm}$ theo đề-xi-mét vuông.
\shortans{486}
\loigiai{
\textbf{Đáp số:} $486$\\
$9 \times 9 \times 6 = 486\text{ dm}^2$.
}
\end{ex}

\begin{ex} %[5H9V51]
Một hình lập phương có diện tích toàn phần là $150\text{ cm}^2$. Tính cạnh của hình lập phương đó theo xăng-ti-mét.
\shortans{5}
\loigiai{
\textbf{Đáp số:} $5$\\
Diện tích 1 mặt: $150 : 6 = 25\text{ cm}^2 \Rightarrow$ Cạnh là $5\text{ cm}$.
}
\end{ex}

\begin{ex} %[5H9V51]
Người ta làm một chiếc hộp hình lập phương không nắp bằng bìa các-tông có cạnh $15\text{ cm}$. Tính diện tích bìa cần dùng theo xăng-ti-mét vuông.
\shortans{1125}
\loigiai{
\textbf{Đáp số:} $1125$\\
$15 \times 15 \times 5 = 1\,125\text{ cm}^2$.
}
\end{ex}

\begin{ex} %[5H9V51]
Một hình lập phương có chu vi một mặt đáy là $28\text{ cm}$. Tính diện tích toàn phần của hình lập phương đó theo xăng-ti-mét vuông.
\shortans{294}
\loigiai{
\textbf{Đáp số:} $294$\\
Cạnh: $28 : 4 = 7\text{ cm}$.\\
$S_{tp} = 7 \times 7 \times 6 = 294\text{ cm}^2$.
}
\end{ex}

\begin{ex} %[5H9C51]
Ghép $8$ khối lập phương nhỏ cạnh $2\text{ cm}$ thành một hình lập phương lớn. Tính diện tích toàn phần của hình lập phương lớn đó theo xăng-ti-mét vuông.
\shortans{96}
\loigiai{
\textbf{Đáp số:} $96$\\
Cạnh hình lập phương lớn: $2 \times 2 = 4\text{ cm}$.\\
$S_{tp} = 4 \times 4 \times 6 = 96\text{ cm}^2$.
}
\end{ex}
